

It is enough to prove that, for every prime ideal \(\mathfrak p\) of
\(B\), the homomorphism
\[
  (A_0)_{\mathfrak p}\otimes_{B_{\mathfrak p}}(A_0)_{\mathfrak p}
  \longrightarrow(A_1)_{\mathfrak p},
\]
with components \(\delta_{0\mathfrak p}\) and
\(\delta_{1\mathfrak p}\), is bijective. Thus we may suppose that
\(B\) is local. \hyperref[proof:sga3-V-4.1-b]{Part~b)} then shows that \((A_0)_{\mathfrak p}\) is
semilocal. Indeed, if \(\mathfrak m\) is one of its maximal ideals,
the others have the form \(\delta_0^{-1}(\mathfrak n)\), where
\(\mathfrak n\) ranges over the prime ideals of \(A_1\) satisfying
\(\delta_1^{-1}(\mathfrak n)=\mathfrak m\). There are at most
\[
  n=[A_1:A_0]
\]
such prime ideals \(\mathfrak n\).

After a faithfully flat base change, we may also suppose that the
residue field of \(B\) is infinite.\footnote{Translator's note: for instance, replace \(B\) by the localization of
\(B[T]\) at the prime ideal generated by the maximal ideal of \(B\), and
\(A_0\), \(A_1\), \(A_2\) by their tensor products with it: this base change is
local and faithfully flat, with infinite residue field \(k(T)\); it preserves
the situation, since forming the equalizer \(B\) of \(\delta_0,\delta_1\)
commutes with flat base change, and the bijectivity to be proved descends
along it. The Gille--Polo re-edition (2011), Exposé~V, refers here to its
Lemma~4.1.2 (see the Addition at the end of e)).}
We can then use \hyperref[lem:sga3-V-4.2]{the following lemma}.

\medskip
\phantomsection\label{lem:sga3-V-4.2}%
\noindent\textbf{Lemma 4.2.}---
\textit{Let \(B\) be a local ring with infinite residue field, let
\(A\) be a semilocal ring, and let \(i:B\to A\) be a homomorphism
which sends the maximal ideal \(\mathfrak n\) of \(B\) into the
radical \(\mathfrak r\) of \(A\). Let \(M\) be a free \(A\)-module of
rank \(n\), and let \(N\) be a \(B\)-submodule of \(M\) which
generates \(M\) as an \(A\)-module. Then \(N\) contains an
\(A\)-basis of \(M\).}

Recall that a sequence \(m_1,\ldots,m_n\) of elements of \(M\) is an
\(A\)-basis if and only if its canonical images in \(M/\mathfrak rM\)
form a basis over \(A/\mathfrak r\). We may therefore replace\footnote{Translator's note: in the last of these replacements the 1970 text has
\(B/\mathfrak m\).}
\[
  M\text{ by }M/\mathfrak rM,\quad
  N\text{ by }N/(N\cap\mathfrak rM),\quad
  A\text{ by }A/\mathfrak r,\quad
  B\text{ by }B/\mathfrak n.
\]
In this case \hyperref[lem:sga3-V-4.2]{the lemma} is elementary. Namely, if
\[
  A=K_1\times\cdots\times K_r,
  \qquad
  M=K_1^n\times\cdots\times K_r^n,
\]
choose, for each \(j\), an element \(x_j\in N\) whose \(j\)-th
component is nonzero. Because the field \(B\) is infinite, a suitable
\(B\)-linear combination \(x\) of the \(x_j\) has all its components
nonzero. Replace \(M\) by \(M/Ax\), and proceed by induction on \(n\).

Apply \hyperref[lem:sga3-V-4.2]{the lemma} with \(A=A_0\), with \(i:B\hookrightarrow A_0\), with
\(M=A_1\) regarded as an \(A_0\)-module through \(\delta_1\), and with
\(N=\delta_0(A_0)\). Indeed, since
\[
  d_0\boxtimes d_1:X_1\longrightarrow X_0\times_YX_0
\]
is a monomorphism, the homomorphism
\[
  A_0\otimes_BA_0\longrightarrow A_1
\]
with components \(\delta_0,\delta_1\) is surjective. This says exactly
that \(\delta_0(A_0)\) generates the \(A_0\)-module \(A_1\).

Choose \(a_1,\ldots,a_n\in A_0\) so that
\(\delta_0(a_1),\ldots,\delta_0(a_n)\) form an \(A_0\)-basis of
\(A_1\). If \(a_1,\ldots,a_n\) form a \(B\)-basis of \(A_0\), then
\(A_0\otimes_BA_0\to A_1\) sends the basis
\[
  (1\otimes a_i)_{1\leq i\leq n}
\]
to the basis
\[
  (\delta_0(a_i))_{1\leq i\leq n},
\]
and is therefore bijective. Thus, if
\[
  \varepsilon:\mathbb Z^n\longrightarrow A_0
\]
is the homomorphism of abelian groups which sends the standard basis
of \(\mathbb Z^n\) to \(a_1,\ldots,a_n\), it remains only to prove
that the map
\[
  B\otimes_{\mathbb Z}\mathbb Z^n\longrightarrow A_0
\]
with components \(i\) and \(\varepsilon\) is bijective.
